Scrivere un programma che, dato un vettore di interi, ne produca due contenenti, rispettivamente, i numeri pari e dispari del vettore iniziale. 

Scrivere una funzione C che, dato in ingresso un vettore di interi di lunghezza N, un numero intero x ed altri parametri se necessario, restituisca 1 se il vettore contiene x, 0 altrimenti.

Scrivere un programma che, dato un vettore (sorg[]) di N double, ne generi un secondo (dest[]) così definito: l'elemento 0 di dest è dato dal prodotto degli elementi 0 e 1 di sorg, l'elemento 1 di dest è dato dal prodotto degli elementi 2 e 3 di sorg, e così via. Per esempio, se sorg[] = {0.5,1.0,1.5,2.0} si deve generare dest[] = {0.5,3.0}.
Si supponga che il numero di elementi di sorg sia pari. Si visualizzi,poi, il vettore dest. 

Scrivere un programma che chieda all’utente due array di 3 elementi e
quindi calcoli la matrice 3 per 3 risultato del prodotto riga per colonna dei
due array.  Il programma dovra’ quindi stampare a video la matrice e la sua
diagonale.